\begin{abstract}
Структура спільного цитування загалом вважається стабільним пошуковим сигналом у правових інформаційних системах.
Ми перевіряємо це припущення лонгітудинально, побудувавши UA-StatuteRetrieval – бенчмарк, що вимірює прогнозованість спільного цитування на 20 щорічних зрізах (2007–2026) із 396 мільйонів посилань на статті кодексів у 101 мільйоні рішень українських судів.
Використовуючи протокол leave-one-out на повному дводольному графі цитувань, ми виявляємо, що MRR за Adamic-Adar знижується на 33\% на \emph{фіксованому} наборі статей (з 0.43 до 0.29) і на 47\% за темпорального розділення train/test (з 0.51 до 0.27) – що підтверджує справжнє темпоральне згасання, а не композиційний зсув чи артефакт оцінювання.
Згасання є нерівномірним: кримінальне процесуальне право зберігає стабільні патерни спільного цитування (MRR $\sim$0.40), тоді як цивільне право деградує з 0.35 до 0.15, що збігається із судовою реформою 2017 року.
Статті-хаби ($>$100K цитувань) протистоять згасанню, але статті середньої частоти (1K–10K) – практичний фронтир пошуку – втрачають половину прогнозованості.
Текстовий базелайн BM25 згасає ще швидше (31\%), а аналіз дрейфу ембедингів за допомогою E5-large виявляє семантичний зсув у 4.3\% у контексті цитування статей, що надає механістичне пояснення спостережуваного згасання.
Бенчмарк опубліковано за адресою \url{https://huggingface.co/datasets/overthelex/ua-statute-retrieval}.
\end{abstract}

\section{Вступ}
\label{sec:intro}

Структура спільного цитування – патерн того, які статті нормативних актів цитуються разом у судових рішеннях – лежить в основі багатьох правових інформаційно-пошукових систем~\citep{tang2024casegnn,ho2025incorporating,barron2025bridging}.
Імпліцитне припущення полягає в тому, що ця структура забезпечує стабільний сигнал: якщо статті A і B часто цитувались разом у минулих рішеннях, вони залишаються релевантними для подібних майбутніх справ.

Ми перевіряємо це припущення емпірично.
З Єдиного державного реєстру судових рішень (ЄДРСР, 101M рішень, 2007–2026) ми витягуємо 396 мільйонів посилань на статті кодексів та будуємо щорічні бенчмарки пошуку, що охоплюють 20 років.
Наш протокол оцінювання – це leave-one-out прогнозування цитувань на повному дводольному графі цитувань: для кожного рішення ми маскуємо кожну цитовану статтю та вимірюємо, чи методи на основі спільного цитування можуть відновити її з решти цитувань.
Хоча це є проксі-задачею для пошуку нормативних актів, а не кінцевою задачею користувача (яка починається з фактів справи, а не часткових цитувань), це безпосередньо вимірює темпоральну стабільність сигналу спільного цитування, що лежить в основі багатьох пошукових систем.

Центральний висновок полягає в тому, що прогнозованість спільного цитування \emph{згасає}: ті самі методи пошуку, застосовані до тих самих статей, дають прогресивно гірші результати з часом.
Це не є тривіальним наслідком зміни складу статей – ми підтверджуємо через абляцію, що згасання зберігається на фіксованому наборі статей (зниження на 33\%) та посилюється за належного розділення train/test (зниження на 47\%).

Згасання концентрується вздовж двох вимірів:
\begin{enumerate}[nosep,leftmargin=*]
  \item \emph{Домен}: кримінальне процесуальне право протистоїть згасанню, тоді як цивільне та адміністративне право деградують швидко.
  \item \emph{Частота}: статті-хаби ($>$100K цитувань) зберігають прогнозованість; статті середньої частоти (1K–10K) – практичний виклик пошуку – втрачають половину сигналу.
\end{enumerate}

Аналіз дрейфу ембедингів за допомогою багатомовної моделі E5-large підтверджує механізм: семантичний контекст, у якому цитуються статті, зсувається на 4.3\% за 12 років, причому цивільне процесуальне право дрейфує найшвидше – що безпосередньо віддзеркалює траєкторії згасання MRR.

Ці висновки мають прямі імплікації для проєктування правових пошукових систем: графові методи не слід розглядати як статичні індекси, а потрібно постійно оновлювати, і необхідні комплементарні текстові методи для хвоста середньої частоти у цивільному/адміністративному праві, де сигнал спільного цитування є найслабшим – хоча як BM25, так і спільне цитування деградують з часом.

\section{Суміжні роботи}
\label{sec:related}

\paragraph{Бенчмарки правового пошуку.}
BSARD~\citep{lotfi2024bsard} надає 1{,}108 запитів для пошуку статей бельгійського законодавства.
SCALE~\citep{rasiah2023scale} та LeSICiN~\citep{paul2021lesicin} оцінюють пошук у швейцарському та індійському контекстах відповідно.
LEXTREME~\citep{niklaus2023lextreme} та LegalBench~\citep{guha2023legalbench} охоплюють класифікацію, але не пошук.
Усі оцінюють в одній часовій точці; жоден не вимірює темпоральну стабільність.

\paragraph{Правові мережі цитувань.}
Fowler et al.~\citep{fowler2007network} досліджують мережі прецедентів Верховного суду США; Coupette et al.~\citep{coupette2021measuring} аналізують еволюцію європейського законодавства; Winkels et al.~\citep{winkels2011determining} вивчають авторитетність нідерландської судової практики.
CaseGNN++~\citep{tang2024casegnn} застосовує графові нейронні мережі до правового пошуку.
Жоден не публікує графи цитувань як бенчмарки пошуку та не досліджує темпоральну деградацію.

\paragraph{Прогнозування зв'язків.}
Ми адаптуємо Common Neighbors та Adamic-Adar~\citep{adamic2003friends} із задачі прогнозування зв'язків у соціальних мережах~\citep{liben2007link} до дводольних правових графів.
Наш внесок полягає в лонгітудинальному застосуванні та відкритті того, що прогнозованість є нестаціонарною.
